\documentclass[12pt]{article}
\usepackage{fullpage}
\usepackage{amsmath,amssymb,amsthm}

\newtheorem{theorem}{Theorem}
\newtheorem{lemma}{Lemma}

\newcommand{\RR}{\mathbb R}
\newcommand{\Vol}{\operatorname{Vol}}
\newcommand{\Dil}{\operatorname{Dil}}
\newcommand{\Mass}{\operatorname{Mass}}

\begin{document}

\begin{center}
{\bf A four-dimensional obstruction to area contraction}
\end{center}

We shall write \(V(R)=R_1R_2R_3R_4\) and
\(V(S)=S_1S_2S_3S_4\).  The letter \(\kappa\) will denote the
small constant to be chosen at the end.

We use the following standard form of Guth's tightening (or
two-parameter sweepout) estimate for rectangles.  A two-parameter
sweepout by relative \(2\)-cycles is a chain-level family in which points
of a parameter rectangle give relative \(2\)-cycles in \(S\), motion in
the first parameter direction through the whole parameter interval gives
relative \(3\)-chains, and the whole parameter rectangle gives a relative
\(4\)-chain.  The degree is the coefficient of the fundamental class in
\(H_4(S,\partial S;\mathbb Z)\) represented by this \(4\)-chain.

\begin{theorem}[Guth's \(4\)-dimensional tightening estimate]\label{T}
There is a dimensional constant \(c_0>0\) with the following property.
Let \(S\subset \RR^4\) be a rectangle with sides
\(S_1\le S_2\le S_3\le S_4\), and let \(C\) be a two-parameter sweepout of
degree \(1\).  Let \(M_2\) be an upper bound for the masses of all
relative \(2\)-cycles in the sweepout, let \(M_3\) be an upper bound for
the masses of the relative \(3\)-chains swept out by moving once across
the first parameter direction, and let \(M_4\) be the mass of the total
relative \(4\)-chain.  Then
\[
 \frac{M_2}{S_3S_4}+
 \frac{M_3}{S_1S_3S_4}+
 \frac{M_4}{S_1S_2^{1/2}S_3^{3/2}S_4}\ge c_0. \tag{1}
\]
\end{theorem}

This is the \(n=4,k=2\) cup-square/sweepout estimate in Guth's
``area-contracting maps between rectangles'' theorem.  For reference, it
is proved by subdividing the parameter rectangle, replacing the chains on
successive skeleta by Federer--Fleming fillings in the target rectangle,
and using the rectangular isoperimetric profile.  The scale
\((S_2S_3)^{1/2}\) in the tightening of the \(3\)-chains is what produces
the factor \(S_2^{1/2}S_3^{3/2}\) in (1).  The estimate is invariant under
independent rescaling of the four coordinates of \(S\), so the constant is
universal.

We now apply the theorem to the map \(f\).  Give \(R\) coordinates
\((x_1,x_2,x_3,x_4)\) adapted to its sides.  Slice
\[
 R=[0,R_1]\times[0,R_2]\times[0,R_3]\times[0,R_4]
\]
by the first two coordinates.  Thus a parameter point \((a,b)\) gives the
relative \(2\)-chain
\(f(\{a\}\times\{b\}\times[0,R_3]\times[0,R_4])\), moving \(a\) from
\(0\) to \(R_1\) with \(b\) fixed gives the relative \(3\)-chain
\(f([0,R_1]\times\{b\}\times[0,R_3]\times[0,R_4])\), and the total
parameter rectangle gives \(f(R)\).  Since \(f:(R,\partial R)\to
(S,\partial S)\) has degree \(1\), this sweepout has degree \(1\).

Because \(\Dil_2(f)\le1\), every \(2\)-dimensional slice
\(\{(x_1,x_2)\}\times[0,R_3]\times[0,R_4]\) has image mass at most
\(R_3R_4\).  Also \(\Dil_2(f)\le1\) implies \(\Dil_j(f)\le1\) for
\(j=3,4\): at each differentiability point the product of the two largest
singular values is at most \(1\), hence all products of three or four
singular values are at most \(1\).  Therefore
\[
 M_2\le R_3R_4,\qquad
 M_3\le R_1R_3R_4,\qquad
 M_4\le R_1R_2R_3R_4. \tag{2}
\]
Substituting (2) into (1) gives
\[
 \frac{R_3R_4}{S_3S_4}
 +\frac{R_1R_3R_4}{S_1S_3S_4}
 +\frac{R_1R_2R_3R_4}{S_1S_2^{1/2}S_3^{3/2}S_4}
 \ge c_0. \tag{3}
\]

Choose
\[
 \kappa=\min\left\{\frac{1}{10},\,\frac{c_0}{4}\right\}.
\]
Assume \(R_1\le \kappa S_1\).  If the first desired alternative fails,
then \(R_3R_4\le \kappa S_3S_4\).  The first two terms in (3) are then at
most
\[
 \kappa \quad\hbox{and}\quad
 \kappa^2,
\]
respectively.  Hence the last term in (3) is at least
\(c_0-\kappa-\kappa^2\ge c_0/2\).  Since \(\kappa\le c_0/4\), this gives
\[
 R_1R_2R_3R_4
 \ge \frac{c_0}{2}\,S_1S_2^{1/2}S_3^{3/2}S_4
 > \kappa\,S_1S_2^{1/2}S_3^{3/2}S_4 .
\]
Thus, if \(R_1\le\kappa S_1\), either
\[
 R_3R_4>\kappa S_3S_4
\]
or
\[
 R_1R_2R_3R_4>\kappa S_1S_2^{1/2}S_3^{3/2}S_4.
\]
This is the required assertion, with the problem's constant \(k\) equal
to \(\kappa\).

\begin{thebibliography}{9}

\bibitem{GuthThesis}
L. Guth, \emph{Area-contracting maps between rectangles}, Ph.D. thesis,
Massachusetts Institute of Technology, 2005.

\bibitem{GuthWidth}
L. Guth, \emph{The width-volume inequality}, Geom. Funct. Anal. 17
(2007), 1139--1179.

\bibitem{GuthExpanding}
L. Guth, \emph{Area-expanding embeddings of rectangles}, arXiv:0710.0403.

\end{thebibliography}

\end{document}
